\documentclass[12pt,a4paper]{article}
\usepackage{fontawesome5}
\providecommand{\faPencilAlt}{\faIcon{pencil-alt}}
\usepackage{styles/couleurs}
\usepackage{styles/mise_en_page}
\usepackage{styles/boites}
\usepackage{styles/maths}
\usepackage{styles/figures}
\usepackage{enumitem}
\usepackage{array}
\usepackage{tabularx}
\usepackage{amsmath}
\usepackage{amssymb}
\setlist[itemize]{leftmargin=1.4em,itemsep=0.15em,topsep=0.2em}
\setlist[enumerate]{leftmargin=*,itemsep=0.25em,topsep=0.25em}
\begin{document}

\begin{center}
{\sffamily\LARGE\bfseries\color{TitreC}Séquence 4 -- Calcul littéral}\\[0.2em]
{\sffamily\large\bfseries\color{SavoirC}Écrire, simplifier et réduire}
\end{center}

\section{Expressions littérales}

\definition{Une \textbf{expression littérale} est une expression contenant une ou plusieurs lettres.
Une lettre représente un nombre dont la valeur peut varier ou être encore inconnue.}

\exemple{Le programme « multiplier un nombre par \(3\), puis ajouter \(5\) »
se traduit par :
\[
3x+5,
\]
où \(x\) représente le nombre de départ.}

\subsection{Calculer la valeur d'une expression}

\methode{Pour calculer la valeur d'une expression :
\begin{enumerate}
  \item remplacer chaque lettre par la valeur donnée ;
  \item écrire des parenthèses autour d'un nombre négatif ;
  \item effectuer le calcul en respectant les priorités.
\end{enumerate}}

\exemple{Pour \(x=-2\) :
\[
3x+5=3\times(-2)+5=-6+5=-1.
\]}

\section{Conventions d'écriture}

\propriete{Dans une expression littérale :
\[
3\times x=3x,
\qquad
x\times5=5x,
\qquad
1\times x=x,
\]
\[
x\times x=x^2,
\qquad
3\times(x+2)=3(x+2).
\]}

\remarque{On écrit généralement le nombre avant la lettre.

Une addition ne se transforme pas en multiplication :
\[
3+x\neq3x.
\]}

\definition{Les nombres qui multiplient les lettres sont appelés les
\textbf{coefficients}. Dans \(5x\), le coefficient de \(x\) est \(5\).}

\section{Réduire une expression}

\definition{Réduire une expression, c'est regrouper les termes de même nature
afin de l'écrire avec le moins de termes possible.}

\propriete{On peut regrouper :
\[
3x+5x=(3+5)x=8x,
\]
\[
4x^2-2x^2=(4-2)x^2=2x^2,
\]
et les nombres entre eux.

On ne peut pas regrouper \(x\), \(x^2\) et les nombres.}

\methode{Pour réduire :
\begin{enumerate}
  \item repérer les termes en \(x^2\), les termes en \(x\) et les nombres ;
  \item additionner séparément leurs coefficients ;
  \item conserver les catégories qui ne peuvent pas être regroupées.
\end{enumerate}}

\exemple{
\[
\begin{aligned}
7x+3-2x+5
&=(7x-2x)+(3+5)\\
&=5x+8.
\end{aligned}
\]}

\exemple{
\[
\begin{aligned}
5x^2+2x-3x^2+7x
&=(5x^2-3x^2)+(2x+7x)\\
&=2x^2+9x.
\end{aligned}
\]}

\erreur{
\[
3x+5\neq8x
\]
car \(3x\) et \(5\) ne sont pas de même nature.

De même :
\[
2x+3x^2\neq5x^3.
\]}

\section{Développer une expression}

\definition{Développer, c'est transformer un produit en une somme ou une différence.}

\propriete{Pour tous nombres \(k\), \(a\) et \(b\) :
\[
k(a+b)=ka+kb
\]
et
\[
k(a-b)=ka-kb.
\]}

\methode{Pour développer, on multiplie chacun des termes placés dans la parenthèse
par le facteur situé devant la parenthèse.}

\exemple{
\[
3(2x+5)=3\times2x+3\times5=6x+15.
\]}

\exemple{On peut ensuite réduire :
\[
\begin{aligned}
4(x+3)+2x
&=4x+12+2x\\
&=6x+12.
\end{aligned}
\]}

\section{Expressions équivalentes et preuve}

\definition{Deux expressions sont \textbf{équivalentes} lorsqu'elles ont la même valeur
pour toute valeur autorisée de la lettre.}

\exemple{
\[
4(x+3)+2x=6x+12
\]
pour toute valeur de \(x\), car le développement et la réduction transforment
la première expression en la seconde.}

\remarque{Tester quelques valeurs peut aider à conjecturer ou à repérer une erreur,
mais ne prouve pas une égalité valable pour toutes les valeurs.

En revanche, un seul contre-exemple suffit pour montrer qu'une égalité générale est fausse.}

\exemple{L'égalité
\[
3(x+2)=3x+2
\]
est fausse. Pour \(x=0\) :
\[
3(0+2)=6
\qquad\text{mais}\qquad
3\times0+2=2.
\]}

\aretenir{
\begin{itemize}
  \item Une lettre représente un nombre.
  \item Simplifier l'écriture consiste notamment à utiliser \(3x\), \(x^2\) et \(3(x+2)\).
  \item Réduire consiste à regrouper les termes de même nature.
  \item Développer consiste à distribuer la multiplication à chaque terme de la parenthèse.
  \item Une transformation littérale correcte conserve la valeur de l'expression.
\end{itemize}}

\end{document}
